\chapter{Boxes}
Temporary appendix material.

\section{Box gallery}
\begin{definition}[Risk-sensitive policy]
A policy $\pi_\beta$ is risk-sensitive if it maximizes a distorted expectation
$Q_\beta(x,a)$ rather than the plain expectation $\mathbb{E}[Z^\pi(x,a)]$.
\end{definition}

\begin{theorem}[Bellman optimality]
The optimal action-value function $Q^*$ is the unique fixed point of the Bellman
optimality operator $\mathcal{T}$, satisfying $Q^* = \mathcal{T}Q^*$.
\end{theorem}

\begin{proposition}[Monotonicity of the distortion]
If $\beta$ is a non-decreasing distortion function, then the induced risk measure
preserves stochastic dominance between return distributions.
\end{proposition}

\begin{example}[Conditional Value-at-Risk]
Setting $\beta(\tau) = \min(\tau/\eta,\, 1)$ recovers $\mathrm{CVaR}_\eta$, which
averages the worst $\eta$-fraction of outcomes and yields a risk-averse policy.
\end{example}

\begin{remark}
In practice the quantile samples $\tau \sim U([0,1])$ are drawn independently for
the online and target networks, which decorrelates the temporal-difference error.
\end{remark}

\section{Experiment Configuration}
Lorem ipsum...

\section{Additional Tables}
Lorem ipsum...

\section{Experiment Configuration}

Lorem ipsum dolor sit amet, consectetur adipiscing elit. Integer facilisis, erat vel porttitor volutpat, neque erat dignissim ipsum, vitae commodo ligula turpis sed nulla.

\section{Additional Tables}

Lorem ipsum dolor sit amet, consectetur adipiscing elit. Integer facilisis, erat vel porttitor volutpat, neque erat dignissim ipsum, vitae commodo ligula turpis sed nulla.
